\documentclass[conference]{IEEEtran}

% ---------------------------------------------------------------------------
% Packages (all widely available in TeX Live / Overleaf)
% ---------------------------------------------------------------------------
\usepackage{graphicx}
\usepackage{booktabs}
\usepackage{amsmath}
\usepackage{listings}
\usepackage[T1]{fontenc}
\usepackage[utf8]{inputenc}
\usepackage{xcolor}
\usepackage{url}
\usepackage[hidelinks]{hyperref}

\lstset{
  basicstyle=\ttfamily\footnotesize,
  breaklines=true,
  columns=fullflexible,
  keepspaces=true,
  frame=single,
  framesep=4pt,
  xleftmargin=4pt,
}

\begin{document}

\title{EBSearch: Topic-Driven Multi-Video Retrieval and\\Synthesized Reporting over Bilibili}

\author{\IEEEauthorblockN{EnhancedBilibiliSearch Project}
\IEEEauthorblockA{Technical Report\\\today}}

\maketitle

\begin{abstract}
Answering a research question from video usually means watching several videos
and reconciling what they say---slow, and harder still on Bilibili, where the
search API must be cryptographically signed, returns no usable relevance score,
and mixes summarizable explainers with multi-hour course dumps. We present
\emph{EBSearch}, a pipeline that turns a single topic into one cited,
synthesized report. It (i) searches Bilibili through the WBI-signed
\texttt{search/type} endpoint, (ii) ranks and selects the top-$N$ most relevant
and \emph{summarizable} videos with a transparent heuristic (and an optional
cheap-LLM re-rank), (iii) fans out per-video summarization by \emph{reusing} a
deployed AIVideoSummary backend over HTTP---inheriting its cookies, ASR relay,
and content-addressed cache---and (iv) synthesizes the per-video summaries into a
single report with a stronger model, under a strict provenance contract in which
every concrete claim carries a bracketed \texttt{[n]} citation that resolves to a
source video. A key empirical finding, validated live against the production
API, is that Bilibili's \texttt{rank\_score} field is null in practice, so we
rank by \emph{result order} reinforced by hit columns, plays, recency, and a
duration-fit window. The system classifies the topic into one of six report
types (how-to, news, comparison, concept, review, general) and emits a
type-specific schema and rendering. EBSearch shares a unified, stdlib-only
account and metering layer with the backend. We give an architectural and
cost/latency-reasoning evaluation; all numbers are explicitly illustrative
design estimates, not measured benchmarks.
\end{abstract}

\begin{IEEEkeywords}
multi-document summarization, information retrieval, large language models,
citation-grounded generation, Bilibili, report generation.
\end{IEEEkeywords}

% ===========================================================================
\section{Introduction}
% ===========================================================================

A common research task is not ``summarize \emph{this} video'' but ``what do the
videos on \emph{this topic} collectively say?'' Done by hand it means searching,
triaging a noisy result list, watching several candidates, and mentally
reconciling agreements, disagreements, and gaps. EBSearch automates this loop for
Bilibili and emits a single, cited report. The setting imposes concrete
constraints that shape the design.

\begin{itemize}
  \item \textbf{Signed, score-less search.} Bilibili's web search endpoint
  (\texttt{wbi/search/type}) requires a per-day WBI request signature, a
  browser-like \texttt{User-Agent}/\texttt{Referer} to clear risk-control
  (HTTP~412), and login cookies for full access. Crucially, the API's
  \texttt{rank\_score} field is null in practice, so a system cannot rank by it
  and must reconstruct relevance from other signals.
  \item \textbf{Noisy candidate pools.} Raw results interleave concise
  explainers with multi-hour course compilations and off-topic popular clips;
  the former are ideal to summarize, the latter are expensive or irrelevant.
  \item \textbf{Cost and rate limits.} Each summarized video and each strong-LLM
  call costs money or tokens, and free-tier providers impose tokens-per-minute
  caps. The pipeline must be aggressively cost-gated: hard prefiltering, a small
  configurable cap on how many videos are summarized, subtitle-first reuse, and
  \emph{exactly one} strong-model synthesis call.
\end{itemize}

\paragraph{Contributions}
\begin{itemize}
  \item A stdlib-only WBI-signed Bilibili search client and a transparent,
  order-based ranking heuristic that does \emph{not} use the (null)
  \texttt{rank\_score}, with hard prefilters for junk and course-length videos
  (Sections~\ref{sec:search},~\ref{sec:rank}).
  \item A fan-out summarization stage that reuses an existing AIVideoSummary
  backend over HTTP---fetching subtitles itself to avoid platform risk-control,
  submitting jobs in parallel, and polling adaptively so wall-clock approximates
  the slowest single video (Section~\ref{sec:fanout}).
  \item A citation-grounded synthesis stage: one stronger-model call over a
  self-contained evidence pack, with a strict \texttt{[n]}-provenance contract,
  deterministic fallbacks, and adaptive report typing into six type-specific
  schemas (Section~\ref{sec:synth}).
  \item Reuse of the unified, stdlib-only account/metering layer shared with the
  backend, with report-level pre-authorize-then-settle billing
  (Section~\ref{sec:account}).
\end{itemize}

% ===========================================================================
\section{System Architecture}
\label{sec:arch}
% ===========================================================================

EBSearch is a four-stage pipeline (Table~\ref{tab:stages}, Fig.~\ref{fig:flow})
orchestrated by a single \texttt{research(topic, cfg)} entry point and exposed as
a CLI, a REST API, and a standalone web page. The core is standard-library only;
FastAPI/uvicorn are optional server dependencies. Each stage is cost-gated, and
exactly one stage makes a strong-model call.

\begin{table}[t]
\centering
\caption{EBSearch pipeline stages.}
\label{tab:stages}
\footnotesize
\begin{tabular}{@{}llp{3.5cm}@{}}
\toprule
\textbf{Stage} & \textbf{Output} & \textbf{Responsibility} \\
\midrule
Search   & \texttt{VideoHit[]}  & WBI-signed multi-plan search; dedup in order \\
Rank/Select & \texttt{ScoredHit[]} & Hard prefilter + heuristic score (+ optional LLM rerank); top-$N$ \\
Fan-out  & \texttt{VideoSummary[]} & Reuse AIVideoSummary backend per video (parallel submit, adaptive poll) \\
Classify & report type          & Topic $\rightarrow$ one of six types (heuristic + optional LLM) \\
Synthesize & \texttt{TopicReport} & One strong-model, cited, type-specific report \\
Render   & Markdown             & Adaptive presentation by report type \\
\bottomrule
\end{tabular}
\end{table}

\begin{figure}[t]
\centering
\footnotesize
\begin{verbatim}
   topic
     |
  [Search]  WBI-signed search/type  (multi-plan,
     |       duration filter, dedup in result order)
     v
   VideoHit[]  (rank_score is NULL -> use position)
     |
  [Rank/Select]  hard prefilter (junk/short/long/
     |           off-topic) -> heuristic score
     |           (position, hit_columns, play,
     |            recency, duration-fit) -> top-N
     |           [+ optional cheap-LLM rerank]
     v
   ScoredHit[N]
     |
  [Fan-out]  per video: fetch subtitle (WBI) ->
     |        POST /api/jobs to AIVideoSummary backend
     |        (transcript job; force_asr url job if
     |         only AI captions + ASR enabled);
     |        parallel submit, adaptive poll
     v
   VideoSummary[]  (reuses backend cookies/relay/cache)
     |
  [Classify] topic -> {how_to,news,comparison,
     |                  concept,review,general}
     v
  [Synthesize] ONE strong-model call over an
     |          evidence pack; strict [n] citations
     v
   TopicReport ---> [Render] adaptive Markdown
\end{verbatim}
\caption{Topic-to-report data flow. The fan-out stage reuses the AIVideoSummary
backend over HTTP; synthesis is the single strong-model call.}
\label{fig:flow}
\end{figure}

% ===========================================================================
\section{Methodology and Implementation}
\label{sec:method}
% ===========================================================================

\subsection{WBI-signed search}
\label{sec:search}

The search client is standard-library only. It first derives the daily
\emph{mixin key} from \texttt{/x/web-interface/nav}: the basenames of two
returned image URLs are concatenated and reordered by a fixed 64-entry
permutation table, and the first 32 characters become the key (cached per client
instance for the day). Each request is then signed by appending \texttt{wts} (a
timestamp) and \texttt{w\_rid}---the MD5 of the sorted, sanitised,
percent-encoded query string concatenated with the mixin key:

\begin{lstlisting}[language=Python]
mixin_key = "".join(s[i] for i in MIXIN)[:32]
params["wts"] = int(time.time())
q = urlencode(sorted(sanitize(params)))
params["w_rid"] = md5((q + mixin_key)).hexdigest()
\end{lstlisting}

Requests carry a browser \texttt{User-Agent}, the Bilibili \texttt{Referer}, and
(when configured) cookies parsed from a Netscape file, clearing risk-control. A
\emph{search plan} (keyword, order, duration filter, page count) is the unit of
work; the query stage can emit several plans (raw, plus optional native
``suggest'' phrasings and optional LLM query expansion), which the client
executes with throttling and merges into a deduplicated list in
\emph{first-seen} (approximately relevance) order. The observed caps are 1000
results / 50 pages / 20 per page. Rows with an empty \texttt{bvid} or zero
duration are dropped at parse time.

\subsection{Ranking without a relevance score}
\label{sec:rank}

The central empirical finding---validated live, since the canonical community
API documentation had been taken down---is that the search API's
\texttt{rank\_score} is null in practice. EBSearch therefore reconstructs
relevance from the signals that are actually present. Selection has two layers,
both cheap and the second optional.

\paragraph{Hard prefilter} A candidate is rejected if it has no \texttt{bvid},
zero or sub-threshold duration (drops clips/shorts), a duration above the
course cutoff (drops multi-hour compilations that are not cheaply
summarizable), near-zero plays, or is \emph{off-topic}---defined as having
neither a title/tag hit flagged by Bilibili nor any overlap between the
title/tags and the topic's tokens (ASCII words plus CJK runs and bigrams).

\paragraph{Transparent weighted heuristic} Survivors receive a weighted score
$$
S = w_p\,p + w_h\,h + w_v\,v + w_r\,r + w_d\,d,
$$
where the components are: $p$, a position decay $1/(1+0.12\cdot\text{pos})$ that
trusts Bilibili's result order (the strongest honest signal); $h$, a
hit-column strength favouring title/tag matches over description/author; $v$,
log-damped plays normalised to the pool maximum so one viral course cannot
dominate; $r$, an exponential recency decay (an $\sim$18-month half-life for a
fast-moving field); and $d$, a Gaussian duration-fit bell peaking near a typical
explainer length. The default weights are
$(w_p,w_h,w_v,w_r,w_d)=(0.34,0.20,0.16,0.12,0.18)$. Course-length videos that
survive a widened window get a multiplicative penalty so they sink below
explainers. The reasons behind every score are retained for transparency.

\paragraph{Optional cheap-LLM re-rank} When enabled, only the heuristic's top
slice is sent to a cheap model for one structured re-rank call; any error falls
back safely to the heuristic order. This bounds cost (the full pool is never
sent) and never hard-fails.

\subsection{Fan-out summarization by backend reuse}
\label{sec:fanout}

Rather than re-implementing transcription, EBSearch \emph{reuses} a deployed
AIVideoSummary backend over HTTP. Because the two services are co-located, the
fan-out can hit the backend at \texttt{127.0.0.1} and inherit its Bilibili
cookies, its overseas ASR relay, and its content-addressed cache---so
re-running a topic or summarizing a previously-seen video is free.

\paragraph{Subtitle policy} To avoid triggering Bilibili's per-video risk-control
through the backend's heavy yt-dlp download, the fan-out \emph{itself} fetches
each video's subtitles via light WBI player-API calls and classifies them as
\texttt{human}, \texttt{ai}, or \texttt{none}. Human subtitles (best quality,
free) are always used by submitting a \texttt{transcript} job. When ASR is
enabled, videos with only AI captions are re-transcribed by submitting a
\texttt{force\_asr} \texttt{url} job (the backend then skips its subtitle-first
path and runs Whisper). When ASR is disabled, AI captions are used as-is rather
than nothing. The duration filter and selection cap remain the main cost gates.

\paragraph{Parallel submit, adaptive poll} Subtitle fetch and job submission run
concurrently across videos through a capped thread pool (the previous serial,
one-second-per-video loop was pure latency); the pool is capped to stay gentle on
the WBI endpoint. The pipeline then polls outstanding jobs, fast at first (cache
hits finish in seconds) and backing off to a steady cadence, until each job is
done, errors, or hits a per-video timeout. The net effect is that the fan-out
wall-clock is dominated by the \emph{slowest single video} rather than the sum
of all of them. Each result records its transcript origin (\texttt{subtitle} vs
\texttt{asr}).

\subsection{Adaptive report typing}

Before synthesis the topic is classified into one of six report types
(\texttt{how\_to}, \texttt{news}, \texttt{comparison}, \texttt{concept},
\texttt{review}, \texttt{general}) by a transport-light classifier: at most one
cheap structured LLM call, always backed by a pure-Python keyword heuristic whose
rules are ordered by specificity so mixed-signal topics resolve sensibly. The
label space is closed; anything not confidently bucketed is \texttt{general}. The
type drives both the synthesis schema and the rendering.

\subsection{Citation-grounded synthesis}
\label{sec:synth}

Synthesis makes \emph{exactly one} call to a stronger model (e.g.\
DeepSeek-v4-pro, or Groq's \texttt{gpt-oss-120b} via the overseas relay) and is
engineered so that faithfulness is structural rather than hoped-for.

\paragraph{Self-contained evidence pack} The model is handed a compact JSON
\emph{evidence pack} containing only, per video, a 1-based citation number
\texttt{n}, its \texttt{bvid}, title, author, duration, TL;DR, capped key points
and keywords, and chapter titles with timestamps. The prompt forbids any fact
not in the pack.

\paragraph{Provenance contract} Every concrete claim must carry a bracketed
\texttt{[n]} marker that resolves to the evidence pack's video \texttt{n} (and
thus to \texttt{sources[n-1]}); the model is told to emit \texttt{[n]} only, never
raw BV numbers or titles. Timestamps may come only from chapter
\texttt{start}/\texttt{t} values, so the model cannot invent timecodes. The
report separates \emph{consensus} from \emph{disagreement} and is required to
flag minority/dissenting views explicitly. We chose this structured-sections
shape over a comparison-matrix-only or a flowing-brief prompt because it maps
1:1 onto the \texttt{TopicReport} object, forces the consensus/dissent split, and
makes provenance checkable per claim.

\paragraph{Type-specific blocks} The base schema (overview, themes, consensus,
disagreements, per-video highlights with timestamps, a ranked watch list, gaps,
sources) is augmented by one type-specific block: \texttt{steps} for how-to,
\texttt{timeline} for news, a \texttt{comparison} matrix for comparison,
\texttt{glossary} for concept, and \texttt{verdicts} for review. The
\texttt{general} type reproduces the base prompt verbatim.

\paragraph{Cost discipline and graceful degradation} Cost is bounded by sending
trimmed evidence (chapter titles only, capped lists) so a report fits within
free-tier token-per-minute limits, capping the synthesis output, and keeping
\texttt{max\_videos} small by default. The synthesis caller falls back from the
primary provider (e.g.\ Groq, which can return 413 on its free tier) to a
direct DeepSeek call so a report is always produced. If the model returns prose
or broken JSON, a fence-aware extractor and deterministic backfill still produce
a well-formed \texttt{TopicReport} (overview from best-effort text; per-video,
sources, and watch list reconstructed straight from the summaries), so the
pipeline never hard-fails on a flaky completion.

\paragraph{Transport} The LLM client is a tiny stdlib OpenAI-compatible chat
function. When a forward proxy is configured it routes the request through an
explicit HTTPS \texttt{CONNECT} tunnel via \texttt{http.client} with
\texttt{set\_tunnel}---\emph{not} \texttt{urllib}'s \texttt{ProxyHandler}, which
does not reliably tunnel HTTPS and can leak the request to the origin IP,
re-triggering a provider geo-block (HTTP~403). TLS stays end-to-end to the
provider, so the proxy never sees the key or body.

\subsection{Adaptive rendering}

The renderer turns a \texttt{TopicReport} into Markdown ordered as the reader's
journey: overview, themes, consensus vs.\ disagreements, a sub-topic by video
coverage matrix, per-video highlights with clickable \texttt{[mm:ss]} deep links
(\texttt{?t=sec}), a ranked watch list, gaps, and sources. The \texttt{[n]}
provenance tags are rewritten into Markdown links to the corresponding source
video, so every claim remains attributable in the rendered output, and the
type-specific block is rendered in a type-appropriate style.

\subsection{Unified account and metering}
\label{sec:account}

EBSearch embeds the \emph{same} stdlib-only account subsystem as the backend and
points it at the \emph{same} SQLite file and the \emph{same} JWT secret, so a
user authenticated against one service is recognised by the other. Storage uses
WAL mode with short \texttt{IMMEDIATE} transactions and a retry-on-locked loop
(two processes, one file); credit changes go through the same atomic
check-and-decrement ledger; tokens are the same hand-rolled HS256 JWTs; and
phone numbers are salted-SHA-256 hashed for PIPL compliance. Billing is
pre-authorize-then-settle at the \emph{report} level: a report's cost is a base
plus a per-video charge times the number of videos (plus a premium-model
surcharge), pre-authorized when the research job starts and refunded on failure.
Because EBSearch already charges its user for the report, its fan-out calls into
the backend present a static service key and are \emph{not} metered again there,
avoiding double-billing.

% ===========================================================================
\section{Evaluation}
\label{sec:eval}
% ===========================================================================

We present an architectural and cost/latency-reasoning evaluation. Every number
below is an \emph{illustrative design estimate} derived from the architecture,
not a measurement, and is labelled as such.
\textbf{[Insert measured end-to-end timings, token counts, and a screenshot of a
rendered report here once a controlled study is run.]}

\subsection{Cost reasoning}

Per topic, the dominant costs are (i) the per-video summaries and (ii) the single
synthesis call. EBSearch minimises both: the duration filter and hard prefilter
cut the pool before any summary is requested; a small \texttt{max\_videos} (4 by
default) caps the number of summaries; subtitle-first reuse means most summaries
are the cheap caption tier; and synthesis is exactly one call over trimmed
evidence. Table~\ref{tab:cost} sketches the relative spend for a default run.

\begin{table}[t]
\centering
\caption{Illustrative per-topic cost decomposition at \texttt{max\_videos}=4,
subtitle-first (\emph{design estimates}, not measurements).}
\label{tab:cost}
\footnotesize
\begin{tabular}{@{}lll@{}}
\toprule
\textbf{Component} & \textbf{Calls} & \textbf{Relative cost} \\
\midrule
Search (WBI)              & several HTTP & negligible (no LLM) \\
Heuristic rank/select     & 0 LLM        & none \\
Per-video summary ($N$)   & $N$ cheap    & low (caption tier) \\
Synthesis                 & 1 strong     & moderate (the main LLM spend) \\
Repeat of seen video/topic& cache hit    & 0 (backend cache) \\
\bottomrule
\end{tabular}
\end{table}

The optional cheap-LLM re-rank and query expansion add at most one cheap call
each and are off by default; the largest \emph{free} quality lever is the
duration filter, which removes course/compilation noise at zero cost.

\subsection{Latency reasoning}

Because the fan-out submits all jobs in parallel and polls adaptively, the
summarization wall-clock approximates the time of the slowest single video, not
the sum---and cache hits return almost immediately. Synthesis is a single call.
Search and ranking are comparatively cheap. Table~\ref{tab:latency} contrasts
the serial baseline with the parallel design for $N$ videos.

\begin{table}[t]
\centering
\caption{Illustrative fan-out wall-clock for $N$ videos (\emph{design
estimates}; $T_i$ is video $i$'s summarization time).}
\label{tab:latency}
\footnotesize
\begin{tabular}{@{}ll@{}}
\toprule
\textbf{Scheme} & \textbf{Summarization wall-clock} \\
\midrule
Serial submit + poll      & $\approx \sum_i T_i$ \\
Parallel submit + adaptive poll (EBSearch) & $\approx \max_i T_i$ \\
All cached                & $\approx$ seconds (poll cadence) \\
\bottomrule
\end{tabular}
\end{table}

\subsection{Qualitative quality observations}

The structured-sections synthesis with a strict provenance contract was chosen
after comparing report shapes on a real fixture; in that exploration the
structured form captured consensus, dissent, and gaps faithfully while keeping
every claim attributable. The order-based ranking, the off-topic guard, and the
course penalty together keep the selected set on-topic and summarizable.
\textbf{[Insert a rendered example report, the selected-video list with scores,
and a faithfulness spot-check here.]}

% ===========================================================================
\section{Related Work}
\label{sec:related}
% ===========================================================================

EBSearch sits at the intersection of information retrieval, multi-document
summarization, and report generation. The per-video stage builds on abstractive
video summarization over transcripts (and, transitively, on robust speech
recognition of the Whisper family~\cite{whisper} via the reused backend). The
topic-level synthesis follows retrieval-augmented generation~\cite{rag}, in which
a generator is constrained to retrieved evidence; EBSearch's evidence pack and
\texttt{[n]} provenance contract are a citation-grounded instance of this idea,
aligned with the broader line of work on report generation that requires every
claim to carry a verifiable source. Relative to closed-source Bilibili search and
summary tools, EBSearch's distinguishing features are its transparent,
score-less ranking heuristic; its reuse of an existing summarization backend over
HTTP; and its adaptive, type-specific, fully cited report schema.

% ===========================================================================
\section{Limitations and Compliance}
\label{sec:limits}
% ===========================================================================

\paragraph{Geo-relay fragility} Synthesis via an overseas provider depends on a
forward proxy that is a single point of failure; the design mitigates this with a
direct-DeepSeek fallback, but a relay outage still degrades the strong path.

\paragraph{Bilibili ToS and risk-control} Search and subtitle fetching rely on
WBI signing, browser-mimicking headers, and login cookies, and are subject to
Bilibili's terms of service and evolving risk-control (the canonical community
API documentation was taken down, so behaviour must be verified empirically and
access throttled conservatively). The null \texttt{rank\_score} is itself an
example of undocumented behaviour that may change.

\paragraph{AIGC labelling} The synthesized report is AI-generated and should be
labelled as such where required by applicable mainland-China AI-generated-content
rules.

\paragraph{No real payment yet} The shared credits ledger accounts for usage but
there is no integrated payment processor; credits are granted administratively or
via invite codes.

\paragraph{Single-node, in-memory job store} Research jobs live in an in-process
registry on a single worker; a restart clears in-flight jobs. Free-tier provider
token-per-minute limits also bound how many videos can be synthesized in one
report (the default cap reflects this), and larger reports need a paid tier.

% ===========================================================================
\section{Conclusion}
\label{sec:conclusion}
% ===========================================================================

EBSearch demonstrates that turning a topic into a trustworthy, cited report over
Bilibili is achievable by combining careful platform engineering with disciplined
reuse. Signing search requests, declining to trust a null relevance score in
favour of result order and a transparent heuristic, reusing an existing
summarization backend (and its cookies, ASR relay, and cache) over HTTP, and
constraining a single strong-model synthesis call to a self-contained evidence
pack with a strict \texttt{[n]} provenance contract together yield reports whose
every claim is attributable, while keeping cost to essentially one strong call
per topic. Adaptive report typing tailors both schema and presentation to the
question. Sharing a stdlib-only account and metering layer with the backend turns
the pipeline into a meterable service without duplicating infrastructure. Future
work includes measured benchmarks to replace the analytical estimates here, a
durable job queue, and paid-tier scaling beyond the current per-report video cap.

% ===========================================================================
\begin{thebibliography}{9}

\bibitem{whisper}
A.~Radford, J.~W.~Kim, T.~Xu, G.~Brockman, C.~McLeavey, and I.~Sutskever,
``Robust Speech Recognition via Large-Scale Weak Supervision,'' OpenAI, 2022.
\url{https://cdn.openai.com/papers/whisper.pdf}

\bibitem{rag}
P.~Lewis et al.,
``Retrieval-Augmented Generation for Knowledge-Intensive NLP Tasks,''
in \emph{Advances in Neural Information Processing Systems (NeurIPS)}, 2020.
\url{https://arxiv.org/abs/2005.11401}

\end{thebibliography}

\end{document}
